\section{Introduction}
\label{sec:introduction}

Research mathematics rarely arrives as an isolated request for a proof.
A researcher must decide which question to pursue, determine what is
already known, interpret the statement correctly, choose representations
that make computation useful, and decide when an approach has stopped
being productive. An experiment that permits these choices studies a
different object from a fixed collection of independently graded proof
problems. Its output is a research process and a collection of artifacts,
and evaluating it requires more than counting final answers.

The Kourovka Notebook~\cite{kourovka21} provides a demanding setting for
such an experiment. It is a long-running collection of problems in group
theory and neighboring areas. Its statements vary substantially in
generality, technical prerequisites, and computational accessibility.
Some have separately posed subparts; some have already been resolved
without an immediately obvious indication in the statement; and some
admit useful partial results without a complete answer. The supplied
September 2026 version of the twenty-first edition was therefore treated
as a source to interpret and investigate, rather than as a benchmark
with a ready-made binary answer key.

The user asked an agent to solve as many previously unsolved Notebook
problems as possible within 48 hours, using a supplied GAP installation,
generally staying within 20 CPU cores and 100\,GB of memory, and keeping
detailed plans, reports, and frequent local Git commits. The agent was
allowed to choose problems, consult the literature, write programs, and
change direction. The research window ran from 10 September 2026,
20:56:46 UTC, to 12 September 2026, 20:56:46 UTC. The final archive was
committed 16 seconds later, after its administrative integrity check.
The manuscript is a separate, later activity.

At the deadline the canonical ledger contained 46 problem entries
labeled complete solution candidates. Several entries cover specified
subparts rather than an entire multipart problem. Multiple covered
subparts of one entry are counted once. The ledger groups \kn{19.63}
with \kn{19.62}; the former's square-closedness criterion was already
published before the experiment and is credited below.
No entry had received outside
review, and priority was unestablished. We retain these qualifications
throughout: a historically recorded candidate is an observable outcome
of the experiment, whereas a correct and new mathematical theorem
requires further judgments.

Three aspects of the experiment are of particular interest. First,
the artifact collection contains explicit mathematical constructions
whose correctness can be considered separately from the agent that
produced them. Second, the session and Git history expose practical
interactions between proof development, computation, and correction.
Third, the experiment produced substantial work that a binary
solved/unsolved classification would obscure: smaller examples of
known phenomena, alternative deductions, sharp bounds in restricted
classes, and complete searches within clearly stated finite ranges.

\subsection{Relation to other experimental settings}

FunSearch combines language-model proposals with executable evaluation
to search for programs producing improved mathematical
constructions~\cite{funsearch}. Formal proof systems provide a different
kind of feedback: AlphaProof, for example, develops proofs within a
formal environment whose accepted derivations can be checked by a
proof assistant~\cite{alphaproof}. The First Proof projects instead use
research-level problems and emphasize the assessment of mathematical
solutions, with the second batch publishing methodology, solutions,
and referee materials~\cite{firstproof,firstproof2}.

The present experiment is an observational case study of a broad
portfolio search. It had neither a fixed executable evaluator for all
targets nor a formal proof kernel accepting the resulting prose
arguments. The agent could also read public sources, including results
relevant to the questions it selected. Consequently, the experiment
does not estimate a held-out success rate, establish an advantage over
another model or a human baseline, or isolate the effect of any
individual tool. Its value lies in the mathematical material and in
the inspectable record of how that material was produced and checked.

\subsection{Reading the paper}

Section~\ref{sec:methods} describes the setting and the limits of the
retained evidence. Section~\ref{sec:trajectory} reconstructs the workflow
and selected changes of direction. Section~\ref{sec:examples} gives
complete representative arguments, including a definability
counterexample and two finite group constructions.
The appendices organize the wider candidate portfolio, partial results,
prior-result deductions, and reproducibility material.
Section~\ref{sec:discussion} discusses what this single run can support
and what a more controlled follow-up should measure.
